\documentclass[11pt]{article}
\usepackage[a4paper,margin=27mm]{geometry}
\usepackage{amsmath,amssymb,amsthm,mathtools,xcolor,booktabs,array}
\newcommand{\PROVED}{\textcolor{green!40!black}{\textbf{PROVED}}}
\newcommand{\OPEN}{\textcolor{red!70!black}{\textbf{OPEN}}}
\newcommand{\AUDIT}{\textcolor{orange!70!black}{\textbf{AUDIT}}}
\newtheorem{theorem}{Theorem}
\newtheorem{proposition}{Proposition}
\newtheorem{definition}{Definition}
\newtheorem{corollary}{Corollary}
\title{The uniform rational-leakage Gamma--Tate theorem (v124)}
\date{14 August 2026}

\begin{document}
\maketitle

\section{The arithmetic leakage operator}

Fix an integer $N\ge2$ and put
\begin{equation}
 T_N=\frac12\log N,
 \qquad
 \delta_N=\frac12\log\frac{N+1}{N},
 \qquad
 E_N=L^2(0,\delta_N).                                      \tag{1}
\end{equation}
Every element of $E_N$ is extended by zero to the real line.  Write
$\Lambda_k=\Lambda^{*k}$ for the $k$-fold Dirichlet convolution of the
von Mangoldt function.  For relatively prime positive integers $a,b$ with
$b>1$ define
\begin{equation}
 \ell_{N,k}(a,b)
 :=\frac1{\sqrt{ab}(\log N)^k}
 \sum_{d\le N/\max(a,b)}
       \frac{\Lambda(bd)\Lambda_k(ad)}d.                    \tag{2}
\end{equation}
The sum is empty when $\max(a,b)>N$.  Hence only the finite set
\begin{equation}
 \mathscr R_N^{\rm leak}
 :=\{a/b:(a,b)=1,\ b>1,\ \max(a,b)\le N\}                 \tag{3}
\end{equation}
can occur.  Let $\tau_u$ denote translation on the zero-extension line,
$(\tau_ue)(t)=e(t-u)$.  Define
\begin{equation}
 L_{N,k}e
 :=\sum_{a/b\in\mathscr R_N^{\rm leak}}
       \ell_{N,k}(a,b)\,\tau_{\log(a/b)}e,                 \tag{4}
\end{equation}
and give different word depths orthogonal output copies:
\begin{equation}
 \mathbf L_Ne:=\bigoplus_{k\ge1}L_{N,k}e.
                                                                    \tag{5}
\end{equation}
For fixed $N$ only finitely many $k$ occur, because
$\Lambda_k(m)=0$ for $k>\log m/\log2$.

For $e\in E_N$ put
\begin{equation}
 \kappa_e(u):=\langle\tau_ue,e\rangle_{L^2(\mathbb R)}.    \tag{6}
\end{equation}

\begin{proposition}[Exact leakage form --- \PROVED]
For every $e\in E_N$,
\begin{align}
 \|\mathbf L_Ne\|^2
 ={}&\sum_{k\ge1}
 \sum_{a/b,c/d\in\mathscr R_N^{\rm leak}}
 \ell_{N,k}(a,b)\overline{\ell_{N,k}(c,d)}
 \kappa_e\!\left(\log\frac{ad}{bc}\right).               \tag{7}
\end{align}
Thus (7) contains all coherent collisions between equal or nearby reduced
fractions.  No orthogonality between overlapping rational cells is being
assumed.
\end{proposition}

\begin{proof}
Expand the squared norm of (4).  The inner product of the two translates
labelled by $a/b$ and $c/d$ is (6) at
$\log((a/b)/(c/d))=\log(ad/bc)$.  Summing the mutually orthogonal word
depths gives (7).  This is exactly the aggregation formula of v123.
\end{proof}

\section{The unreset Gamma form and the two Tate terminals}

Put
\begin{equation}
 a_j=2j+\frac12\qquad(j\ge0).                               \tag{8}
\end{equation}
For $f\in C_c^\infty(\mathbb R)$ let $z_{j,f}$ be the unique causal state
\begin{equation}
 z_{j,f}'=-a_jz_{j,f}+\sqrt{2a_j}\,f,
 \qquad z_{j,f}(-\infty)=0.                                \tag{9}
\end{equation}
Define the complete unreset Gamma energy
\begin{align}
 \mathcal G_\Gamma[f]
 &: =\sum_{j\ge0}\frac1{a_j^2}
       \int_{\mathbb R}|z_{j,f}'(t)|^2\,dt                 \tag{10}\\
 &=\int_{\mathbb R}g_\Gamma(\tau)|\widehat f(\tau)|^2
       \frac{d\tau}{2\pi},\qquad
 g_\Gamma(\tau)=
 \sum_{j\ge0}\frac{2\tau^2}{a_j(a_j^2+\tau^2)}.          \tag{11}
\end{align}
Equation (9) is solved once on the whole line.  In particular, the state
leaving one rational cell is the state entering the next one; (10) does
not reset any state at an internal endpoint.

For a compactly supported $f$ define its two Tate moments
\begin{equation}
 M_+(f)=\int_{\mathbb R}e^{t/2}f(t)\,dt,
 \qquad
 M_-(f)=\int_{\mathbb R}e^{-t/2}f(t)\,dt.                  \tag{12}
\end{equation}
If $I=(L,R)$ contains $\operatorname{supp}f$, the Gram of the two Tate
vectors on $I$ is
\begin{equation}
 G_I=
 \begin{pmatrix}
 e^R-e^L&R-L\\
 R-L&e^{-L}-e^{-R}
 \end{pmatrix}.                                             \tag{13}
\end{equation}
It is strictly positive for $L<R$.  Set
\begin{equation}
 \mathcal T_I[f]
 :=
 \begin{pmatrix}\overline{M_+(f)}&\overline{M_-(f)}\end{pmatrix}
 G_I^{-1}
 \binom{M_+(f)}{M_-(f)}.                                   \tag{14}
\end{equation}
This is the squared norm of the orthogonal Tate component of $f$ in
$L^2(I)$.  The interval $I$ must therefore be fixed before the estimate;
one may not replace it cell by cell or reset the two terminal coordinates
at internal endpoints.

Fix one open interval containing the supports of all $L_{N,k}e$ when
$e\in C_c^\infty(0,\delta_N)$.  Uniformly in the arithmetic of $N$, take
\begin{equation}
 I_N=\left(-\log N,\ \log\frac N2+\delta_N\right),          \tag{15}
\end{equation}
which contains every interval
$\log(a/b)+(0,\delta_N)$ occurring in (4).  Define the Gamma--Tate form
on the physical leakage column by
\begin{equation}
 \mathcal E_{\Gamma T,N}(e)
 :=\sum_{k\ge1}
 \left\{\mathcal G_\Gamma[L_{N,k}e]
             +\mathcal T_{I_N}[L_{N,k}e]\right\}.           \tag{16}
\end{equation}
Its natural form domain is
\begin{equation}
 \mathscr D_N:=\{e\in E_N:\mathcal E_{\Gamma T,N}(e)<\infty\}.
                                                                    \tag{17}
\end{equation}
The space $C_c^\infty(0,\delta_N)$ is contained in $\mathscr D_N$ and is
dense in $E_N$.  Both sides of (16) are constructed from $\Lambda$,
$\Lambda_k$, translations, the Gamma parameters (8), and the two Tate
moments.  No zero of $\xi$ and no inverse of the desired row-(d) operator
occurs.

\section{Uniform rational-leakage Gamma--Tate domination}

\begin{theorem}[URGT$_1$, exact unit form --- \OPEN]
For every integer $N\ge2$ and every $e\in\mathscr D_N$,
\begin{equation}
 \boxed{
        \|\mathbf L_Ne\|^2
        \le \mathcal E_{\Gamma T,N}(e).}                  \tag{18}
\end{equation}
Equivalently, in the order of nonnegative closed quadratic forms on
$E_N$,
\begin{equation}
 \boxed{
        \mathbf L_N^*\mathbf L_N
        \le \mathbf Q_{\Gamma T,N},}                       \tag{19}
\end{equation}
where $\mathbf Q_{\Gamma T,N}$ is the positive self-adjoint operator
associated by the representation theorem to the closed form (16).
The constant in (18)--(19) is exactly one and is independent of
$N$, of the word depth, of the reduced rational labels, and of the cell
profile $e$.
\end{theorem}

The theorem is deliberately stated on the aggregated physical output
(4), not term by term in $(a,b)$.  Therefore it controls destructive and
constructive interference alike.  It is also an all-depth assertion: the
sum over $k$ is part of the operator, not an asymptotic remainder.

\begin{proposition}[Completely expanded scalar form --- \PROVED]
The statement URGT$_1$ is equivalent to the following single explicit
inequality, valid for every $N\ge2$ and
$e\in C_c^\infty(0,\delta_N)$:
\begin{align}
 &\sum_{k\ge1}
 \sum_{a/b,c/d\in\mathscr R_N^{\rm leak}}
 \ell_{N,k}(a,b)\overline{\ell_{N,k}(c,d)}
 \kappa_e\!\left(\log\frac{ad}{bc}\right)                \notag\\
 &\quad\le
 \sum_{k\ge1}\int_{\mathbb R}g_\Gamma(\tau)
     \left|
       \widehat e(\tau)
       \sum_{a/b\in\mathscr R_N^{\rm leak}}
          \ell_{N,k}(a,b)e^{-i\tau\log(a/b)}
     \right|^2\frac{d\tau}{2\pi}                         \notag\\
 &\qquad+
 \sum_{k\ge1}
 \begin{pmatrix}\overline{M_+(L_{N,k}e)}&
                 \overline{M_-(L_{N,k}e)}\end{pmatrix}
 G_{I_N}^{-1}
 \binom{M_+(L_{N,k}e)}{M_-(L_{N,k}e)}.                    \tag{20}
\end{align}
Closure of the two sides extends (20) from the test core to
$\mathscr D_N$.
\end{proposition}

\begin{proof}
The left side is (7).  Translation gives
\[
 \widehat{L_{N,k}e}(\tau)
 =\widehat e(\tau)
   \sum_{a/b}\ell_{N,k}(a,b)e^{-i\tau\log(a/b)}.
\]
Substitution in (11), followed by (12)--(16), gives the right side of
(20).  Thus (18) and (20) agree on the dense test core.  The standard
closure theorem for ordered nonnegative forms gives the equivalence with
(19) and the extension to (17).
\end{proof}

\section{What proving URGT$_1$ would close}

\begin{proposition}[Conditional logical cascade --- \PROVED]
Assume URGT$_1$.  Assume also the already audited identifications of
v117--v123: the G39 word column is the integral-label part of the physical
Euler boundary current; (10) is the unreset Gamma resistance; and (14) is
the conservative moving-Tate terminal.  Then the joint
Gamma--Euler--Tate source map is contractive with the unit Schur
normalization.  Consequently the shorted inverse-metric domination of
v119--v121 holds on every arithmetic threshold cell.
\end{proposition}

\begin{proof}
The physical Euler output decomposes orthogonally into its integral-label
part and its nonintegral rational part.  The integral part is the G39 Witt
current by v119.  The nonintegral part is exactly $\mathbf L_Ne$ by
(2)--(5).  URGT$_1$ charges its complete Gram to the positive unreset
Gamma--Tate form with coefficient one.  The moving-Tate column is
isometric by v117, so no additional norm is created at the two exterior
ports.  Orthogonal addition of the integral and leakage ports therefore
gives the joint unit source contraction.  Shorting a passive source
network preserves passivity, yielding the inverse-metric domination.
\end{proof}

The proposition fixes the exact role of (18).  It does not mark the final
consequences as proved before (18) itself is proved.  In particular, the
status is
\[
 \text{URGT$_1$}
 \Longrightarrow \text{shorted domination}
 \Longrightarrow (8c)
 \Longrightarrow G40
 \Longrightarrow G
 \Longrightarrow \text{row (d)}.
\]

\begin{center}
\begin{tabular}{>{\raggedright\arraybackslash}p{.57\linewidth}
                >{\centering\arraybackslash}p{.14\linewidth}
                >{\raggedright\arraybackslash}p{.19\linewidth}}
\toprule
Statement & Status & Location \\
\midrule
Exact aggregated leakage operator & \PROVED & (2)--(7) \\
Unreset complete Gamma energy & \PROVED & (8)--(11), v121 \\
Two-dimensional Tate terminal form & \PROVED & (12)--(16), v117--v122 \\
Equivalence of operator and expanded forms & \PROVED & (18)--(20) \\
Uniform unit domination URGT$_1$ & \OPEN & (18)/(20) \\
Shorted domination, (8c), $G40$, $G$, row (d) & \OPEN & after URGT$_1$ \\
\bottomrule
\end{tabular}
\end{center}

\paragraph{Audit warning.}
It would be circular to define $\mathbf Q_{\Gamma T,N}$ as
$\mathbf L_N^*\mathbf L_N$ or to insert the desired positive Schur
complement into (16).  Equations (8)--(16) define the right side before
URGT$_1$ and entirely from source-labelled Gamma and Tate data.  The next
step is a proof of (20), not a further reformulation of it.

\end{document}
